\documentclass[11pt]{article}

\usepackage{amsmath,amsthm,amssymb}
\usepackage[margin=1in]{geometry}

\newtheorem{theorem}{Theorem}
\newtheorem{proposition}[theorem]{Proposition}
\newtheorem{corollary}[theorem]{Corollary}
\newtheorem{lemma}[theorem]{Lemma}
\theoremstyle{definition}
\newtheorem{definition}[theorem]{Definition}
\newtheorem{example}[theorem]{Example}
\newtheorem{remark}[theorem]{Remark}
\newtheorem{hypo}{Hypothesis}
\renewcommand{\thehypo}{(H)}

\newcommand{\Hom}{\operatorname{Hom}}
\newcommand{\End}{\operatorname{End}}
\newcommand{\Aut}{\operatorname{Aut}}
\newcommand{\id}{\mathrm{id}}
\newcommand{\C}{\mathcal{C}}
\newcommand{\D}{\mathcal{D}}
\newcommand{\K}{\mathcal{K}}
\newcommand{\E}{\mathcal{E}}
\newcommand{\Sk}{\mathrm{Sk}}
\newcommand{\Set}{\mathbf{Set}}
\newcommand{\Ab}{\mathbf{Ab}}
\newcommand{\ob}{\operatorname{ob}}
\newcommand{\dom}{\operatorname{dom}}
\newcommand{\cod}{\operatorname{cod}}
\newcommand{\orb}{\operatorname{orb}}
\newcommand{\flr}[1]{\lfloor #1 \rfloor}
\newcommand{\1}{\mathbf{1}}

\title{The global-closure obstruction (G) is a cohomology class:\\
the Zappa--Sz\'ep holonomy as $H^2(\Sk_\C;\D)$}
\author{MacBeth\\\small Kodamai}
\date{11 June 2026}

\begin{document}
\maketitle

\begin{abstract}
We resolve the next layer of SEED-Q2. The two-level criterion for a strict
factorization $\K=\C\bowtie\D$ over a wide subcategory $\D$ splits into a fibrewise
freeness condition \textnormal{(L)} and a global closure condition
\textnormal{(G)}; the prior notes proved \textnormal{(L)}$\wedge$\textnormal{(G)}
decides the question but described \textnormal{(G)} only as ``an obstruction.''
Here we \emph{classify} it. Assuming \textnormal{(L)}, every transversal (choice of
free bases) carries a \emph{defect $2$-cochain} $\omega$ valued in $\D$, with
\textnormal{(G)} holding iff $\omega\equiv\id$. We isolate a clean hypothesis
\ref{hyp} --- abelian vertex groups and a \emph{trivial left action} --- under
which the free $\D$-orbits form a category $\Sk_\C$, the vertex automorphism groups
form a presheaf $\D\colon\Sk_\C^{\mathrm{op}}\to\Ab$, the defect is a normalized
$2$-cocycle, and re-choosing the transversal changes it by a coboundary. Hence
$[\omega]\in H^2(\Sk_\C;\D)$ is a transversal-independent invariant and
\[
   \textbf{(G) holds }\iff\ [\omega]=0 .
\]
The set of strict factorizations, when nonempty, is a torsor under
$Z^1(\Sk_\C;\D)$, and modulo inner relabelings under $H^1(\Sk_\C;\D)$ --- the
$H^1$/$H^2$ pattern the conjecture predicted. We compute the ten-morphism rigid
twist explicitly: $\Sk_\C$ is the branch $a\to x\rightrightarrows y$, $\D$ is
$\mathbb Z/2$ concentrated at $a$, $H^2(\Sk_\C;\D)\cong\mathbb Z/2$, and the defect
represents the \emph{nonzero} element; the cross-arrow enlargement of the pairwise
counterexample gives the same class. Finally we delimit honestly: a nontrivial
left action (which already occurs in two-morphism examples) breaks the descent to
$\Sk_\C$, and the obstruction's natural home becomes the \emph{nonabelian} $H^2$ of
the matched pair --- the Kac/Rosebrugh--Wood regime, which we state as the open
continuation. All claims are machine-checked
(\texttt{projects/scratch/cohomology\_holonomy.py}).
\end{abstract}

\section{Setup}

Throughout, $\K$ is a small category and $\D\subseteq\K$ a wide subcategory (the
prescribed right factor). We use the pairwise criterion of the companion note
\cite{pairwise} freely: $\K=\C\bowtie\D$ for some wide $\C$ iff
\begin{itemize}
\item[\textnormal{(L)}] each hom-presheaf $\Hom_\K(-,b)$ is free as a right
$\D$-set; and
\item[\textnormal{(G)}] free bases $B_b\subseteq\Hom_\K(-,b)$ can be chosen so that
$\C=\bigcup_b B_b$ is a wide subcategory.
\end{itemize}
We \emph{assume \textnormal{(L)} throughout}. Our object of study is the residual
condition \textnormal{(G)}.

Recall the structure of a free $\D$-set. The right $\D$-action on $\Hom_\K(-,b)$ is
precomposition, $f\cdot d=f\circ d$ for $d\in\D$. An \emph{orbit} is a connected
component of the action; \textnormal{(L)} says each orbit is free. A morphism
$c\colon a\to b$ \emph{generates} its orbit $o$ iff $d\mapsto c\circ d$ is a
bijection from the $\D$-arrows into $a$ onto $o$; equivalently every $f$ in $o$
factors uniquely as $f=c\circ d$, $d\in\D$. We call $a=\dom(c)$ the \emph{apex} of
$o$ (well defined up to $\D$-isomorphism). Two generators of $o$ differ by a
$\D$-isomorphism on the right: $c'=c\circ\varphi$ with $\varphi$ a $\D$-iso into
the apex. We write $\orb(f)$ for the orbit of $f$.

\begin{definition}[Vertex groups]
For an object $a$ put $G_a:=\Aut_\D(a)$, the group of $\D$-automorphisms of $a$.
When $\D$ is a disjoint union of vertex groups (the single-category case
$\D=\E$ with invertible endomorphisms), $G_a=\End_\D(a)$ and apex $=$ source.
\end{definition}

\begin{definition}[Transversal]
A \emph{transversal} $T$ is a choice of one generator $c_o\in o$ for every orbit
$o$ (over all targets $b$), subject to the \emph{unit normalization}: at the orbit
of $\id_b$ (which is the regular orbit $\D(-,b)$, generated by the $\D$-automorphisms
of $b$) we choose $c=\id_b$. The associated candidate left factor is
$\C_T:=\{c_o\}_o$, the set of chosen generators.
\end{definition}

By \cite[Thm.~1]{pairwise}, $\C_T$ is the left factor of a strict factorization
system $(\C_T,\D)$ iff $\C_T$ is closed under composition; this is exactly
\textnormal{(G)} for the choice $T$.

\section{The defect $2$-cochain (general)}

The following is unconditional (it needs only \textnormal{(L)} and a transversal).

\begin{definition}[Defect]\label{def:defect}
Fix a transversal $T$. For composable generators $c_1\colon a\to b$ and
$c_2\colon b\to e$ in $\C_T$, the composite $c_2\circ c_1\colon a\to e$ lies in some
orbit at $e$; let $\flr{c_2c_1}_T\in T$ be that orbit's chosen generator. By
freeness there is a \emph{unique} $\omega_T(c_2,c_1)\in\D$ with
\begin{equation}\label{eq:defect}
   c_2\circ c_1 \;=\; \flr{c_2 c_1}_T \circ\, \omega_T(c_2,c_1),
   \qquad \omega_T(c_2,c_1)\colon a\to \dom\flr{c_2c_1}_T .
\end{equation}
We call $\omega_T$ the \emph{defect $2$-cochain} of $T$.
\end{definition}

\begin{proposition}[Closure $=$ vanishing of the defect]\label{prop:closure}
\textnormal{(G)} holds for the transversal $T$ (i.e.\ $\C_T$ is a wide subcategory)
if and only if $\omega_T(c_2,c_1)$ is an identity for every composable pair.
\end{proposition}

\begin{proof}
$\C_T$ contains all identities (unit normalization) and is closed under composition
iff $c_2\circ c_1\in\C_T$ for all composable generators. By \eqref{eq:defect},
$c_2\circ c_1=\flr{c_2c_1}_T\circ\omega_T(c_2,c_1)$ with $\flr{c_2c_1}_T\in\C_T$;
since a generator together with a nonidentity $\D$-arrow gives, by uniqueness of
factorization, an element \emph{outside} the generating set unless the $\D$-arrow
is an identity, we have $c_2\circ c_1\in\C_T\iff\omega_T(c_2,c_1)=\id$. (If
$\omega_T(c_2,c_1)=\id$ then $c_2c_1=\flr{c_2c_1}_T\in\C_T$; conversely if
$c_2c_1\in\C_T$ it is a chosen generator, hence \emph{the} generator
$\flr{c_2c_1}_T$ of its orbit, forcing $\omega_T(c_2,c_1)=\id$ by uniqueness.)
\end{proof}

\begin{definition}[Gauge]\label{def:gauge}
The \emph{gauge group} is $\mathcal G:=\prod_{o}\Gamma_o$, where $\Gamma_o$ is the
group of $\D$-automorphisms of the apex of $o$ acting on generators of $o$ by right
multiplication $c_o\mapsto c_o\circ h_o$, and $h$ is normalized ($h_o=\id$ on unit
orbits). $\mathcal G$ acts \emph{freely and transitively} on the set of
transversals: any two transversals differ by a unique gauge (replace each chosen
generator by the other orbit's generator, which differ by a $\D$-iso). Write
$T\cdot h$ for the regauged transversal.
\end{definition}

\section{The clean hypothesis and the cochain complex}

The gauge transformation of $\omega$ involves \emph{post}composition by target
automorphisms --- the left action of the matched pair --- and is nonabelian in
general. We isolate the regime in which it is a coboundary in an abelian complex.

\begin{hypo}\label{hyp}\leavevmode
\begin{enumerate}
\item[(i)] \emph{(abelian vertex groups)} $\D$ is a disjoint union of vertex groups,
$\D(a,a)=G_a=\End_\D(a)$ a group and $\D(a,b)=\varnothing$ for $a\neq b$, with every
$G_a$ abelian; equivalently $\K$ is a single small category with invertible
endomorphisms and $\D=\E$ \cite{cocycle};
\item[(ii)] \emph{(trivial left action)} for every generator $c\colon a\to b$ and
every $\psi\in G_b$, the morphism $\psi\circ c$ lies in $\orb(c)$ --- equivalently
$\psi\circ c=c\circ\varphi_c(\psi)$ for a (necessarily unique) $\varphi_c(\psi)\in
G_a$.
\end{enumerate}
\end{hypo}

Under \ref{hyp}(i) the apex of every orbit equals the common source of its elements
(the $\D$-arrows into $a$ are exactly $G_a$, all $a\to a$), and for composable
generators $c_1\colon a\to b$, $c_2\colon b\to e$ the composite $c_2\circ c_1$
\emph{generates} its orbit: the orbit is $\{c_2c_1\circ\gamma:\gamma\in G_a\}$ and
$\gamma\mapsto c_2c_1\circ\gamma$ is a bijection by freeness \textnormal{(L)}. Hence
the defect $\omega_T(c_2,c_1)$ of \eqref{eq:defect} is forced to lie in the
\emph{group} $G_a=G_{\dom c_1}$ --- the coefficient we cohomologize.

\begin{lemma}[Descent to the orbit category]\label{lem:skel}
Under \ref{hyp}(ii), the assignment $o_2\ast o_1:=\orb(c_2\circ c_1)$ (any
generators $c_i$ of $o_i$, composable) is well defined and associative, and makes
$\Sk_\C$ --- objects $\ob\K$, morphisms the free $\D$-orbits, identities the unit
orbits --- a category. Under \ref{hyp} the rule $a\mapsto G_a$ on objects and
$\varphi_c\colon G_b\to G_a$ on a generator $c\colon a\to b$ extends to a presheaf
$\D\colon\Sk_\C^{\mathrm{op}}\to\Ab$ (a functor on $\Sk_\C^{\mathrm{op}}$).
\end{lemma}

\begin{proof}
\emph{Well-defined composition.} Replacing $c_1$ by $c_1\circ\varphi$
($\varphi$ a $\D$-iso into the apex) leaves $\orb(c_2\circ c_1\circ\varphi)=
\orb(c_2\circ c_1)$ (right $\D$-translation preserves orbits). Replacing $c_2$ by
$c_2\circ\psi$ ($\psi\in G_b$, $b$ the apex of $o_2$): by \ref{hyp}(ii),
$\psi\circ c_1=c_1\circ\varphi_{c_1}(\psi)$, so
$c_2\circ\psi\circ c_1=c_2\circ c_1\circ\varphi_{c_1}(\psi)$, again the same orbit.
Associativity and unitality of $\ast$ are inherited from $\K$ because
$\orb(-)$ depends only on the underlying $\K$-morphism, which associates.

\emph{Presheaf.} $\varphi_c\colon G_b\to G_a$ is a homomorphism: for
$\psi,\psi'\in G_b$,
$c\circ\varphi_c(\psi'\psi)=(\psi'\psi)\circ c=\psi'\circ(\psi\circ c)
=\psi'\circ c\circ\varphi_c(\psi)=c\circ\varphi_c(\psi')\varphi_c(\psi)$, and
cancel $c$ on the left (it is a generator, hence left-cancellable against
distinct $G_a$-translates). Functoriality $\varphi_{c_2\ast c_1}=\varphi_{c_1}
\circ\varphi_{c_2}$ is the same calculation applied to $c_2\circ c_1$, using
\ref{hyp}(ii) twice; $\varphi_{\id}=\id$. Contravariance is built into the typing
$G_b\to G_a$ along $c\colon a\to b$.
\end{proof}

We use the standard cochain complex of a small category with coefficients in a
presheaf \cite{baueswirsching,roos}. For $n\ge 0$,
\[
   C^n(\Sk_\C;\D)=\!\!\prod_{o_n,\dots,o_1\ \text{composable}}\!\! G_{\dom o_1},
\]
\[
   (\delta\sigma)(o_{n+1},\dots,o_1)=
   \varphi_{o_1}\!\big(\sigma(o_{n+1},\dots,o_2)\big)
   +\sum_{i=1}^{n}(-1)^i\,\sigma(o_{n+1},\dots,o_{i+1}\ast o_i,\dots,o_1)
   +(-1)^{n+1}\sigma(o_{n},\dots,o_1),
\]
the first face dropping the rightmost arrow $o_1$ and acting by its restriction,
the last dropping the leftmost $o_{n+1}$; normalized cochains vanish whenever some
$o_i$ is a unit. Concretely
\begin{align}
(\delta^1 h)(o_2,o_1)&=\varphi_{o_1}(h(o_2))-h(o_2\ast o_1)+h(o_1),\label{eq:d1}\\
(\delta^2\omega)(o_3,o_2,o_1)&=\varphi_{o_1}(\omega(o_3,o_2))-\omega(o_3,o_2\ast o_1)
   +\omega(o_3\ast o_2,o_1)-\omega(o_2,o_1).\label{eq:d2}
\end{align}
$\delta^2\circ\delta^1=0$ (verified directly: the eight terms cancel using
$\varphi_{o_2\ast o_1}=\varphi_{o_1}\circ\varphi_{o_2}$);
$H^n(\Sk_\C;\D):=\ker\delta^{n}/\operatorname{im}\delta^{n-1}$.

\section{Main theorem}

A transversal is the same thing as a section of $\Sk_\C$ over $\K$ (a choice of
generator per orbit); the gauge group of Definition~\ref{def:gauge} is exactly the
group $C^1(\Sk_\C;\D)$ of normalized $1$-cochains, $h\colon o\mapsto h_o\in
G_{\dom o}$.

\begin{theorem}[The closure obstruction is a cohomology class]\label{thm:main}
Assume \textnormal{(L)} and \ref{hyp}. Then:
\begin{enumerate}
\item the defect $\omega_T$ descends to a normalized $2$-cochain on $\Sk_\C$ and is
a $2$-cocycle: $\delta^2\omega_T=0$;
\item regauging changes it by a coboundary: $\omega_{T\cdot h}=\omega_T+\delta^1 h$
for every $h\in C^1(\Sk_\C;\D)$;
\item consequently $[\omega]:=[\omega_T]\in H^2(\Sk_\C;\D)$ is independent of the
transversal $T$, and
\[
   \textnormal{(G) holds}\quad\Longleftrightarrow\quad [\omega]=0 ;
\]
\item if nonempty, the set of strict factorization systems $(\C,\D)$ of $\K$ is a
torsor under $Z^1(\Sk_\C;\D)$; modulo the inner relabelings
$B^0=\operatorname{im}\delta^0$ (automorphisms of $(\K,\D)$ fixing objects) it is a
torsor under $H^1(\Sk_\C;\D)$.
\end{enumerate}
\end{theorem}

\begin{proof}
Throughout, additive notation in the abelian groups $G_a$. By the remark after
\ref{hyp}, the defect $\omega_T(c_2,c_1)$ lands in the group $G_{\dom c_1}$. It
depends only on the orbits $o_i=\orb(c_i)$: changing $c_1\mapsto c_1\gamma$
($\gamma\in G_a$) multiplies both sides of \eqref{eq:defect} on the right by
$\gamma$ and rescales the defect accordingly; changing $c_2\mapsto c_2\psi$
($\psi\in G_b$) uses \ref{hyp}(ii), $\psi\circ c_1=c_1\circ\varphi_{c_1}(\psi)$, to
absorb the change into the right $G_a$-action. So $\omega_T$ is a normalized
$2$-cochain on $\Sk_\C$.

\emph{(1) Cocycle.} Factor $c_3\circ c_2\circ c_1$ two ways and use uniqueness.
Writing $\flr{\cdot}$ for $\flr{\cdot}_T$ and $\omega$ for $\omega_T$:
\begin{align*}
c_3\circ(c_2\circ c_1)
&=c_3\circ\big(\flr{c_2c_1}\,\omega(c_2,c_1)\big)
 =\big(c_3\circ\flr{c_2c_1}\big)\,\omega(c_2,c_1)\\
&=\flr{c_3\,\flr{c_2c_1}}\,\omega\!\big(c_3,\flr{c_2c_1}\big)\,\omega(c_2,c_1),
\end{align*}
while
\begin{align*}
(c_3\circ c_2)\circ c_1
&=\big(\flr{c_3c_2}\,\omega(c_3,c_2)\big)\circ c_1
 =\flr{c_3c_2}\circ\big(\omega(c_3,c_2)\circ c_1\big)\\
&=\flr{c_3c_2}\circ c_1\circ\varphi_{c_1}\!\big(\omega(c_3,c_2)\big)
 =\flr{\flr{c_3c_2}\,c_1}\,\omega\!\big(\flr{c_3c_2},c_1\big)\,
   \varphi_{c_1}\!\big(\omega(c_3,c_2)\big),
\end{align*}
the third equality by \ref{hyp}(ii). The two leading generators agree
($\flr{c_3\flr{c_2c_1}}=\flr{c_3c_2c_1}=\flr{\flr{c_3c_2}c_1}$ since $\ast$ is
associative in $\Sk_\C$). By uniqueness of the $\D$-part,
\[
   \omega\!\big(c_3,\flr{c_2c_1}\big)+\omega(c_2,c_1)
   =\omega\!\big(\flr{c_3c_2},c_1\big)+\varphi_{c_1}\!\big(\omega(c_3,c_2)\big),
\]
which is exactly $(\delta^2\omega)(o_3,o_2,o_1)=0$ in \eqref{eq:d2}.

\emph{(2) Coboundary.} Regauge $c_i\mapsto c_i\,h_{o_i}$, $h_{o_i}\in G_{\dom c_i}$.
Then $\flr{\cdot}\mapsto\flr{\cdot}\,h_{o_2\ast o_1}$ and
\begin{align*}
(c_2 h_{o_2})\circ(c_1 h_{o_1})
&=c_2\circ(h_{o_2}\circ c_1)\circ h_{o_1}
 =c_2\circ c_1\circ\varphi_{c_1}(h_{o_2})\,h_{o_1}\\
&=\flr{c_2c_1}\,\omega_T(c_2,c_1)\,\varphi_{o_1}(h_{o_2})\,h_{o_1}
 =\big(\flr{c_2c_1}h_{o_2\ast o_1}\big)\,
  \big(h_{o_2\ast o_1}^{-1}\,\omega_T(c_2,c_1)\,\varphi_{o_1}(h_{o_2})\,h_{o_1}\big),
\end{align*}
using \ref{hyp}(ii) for $h_{o_2}\circ c_1$. Reading off the $\D$-part for the
regauged transversal,
\[
   \omega_{T\cdot h}(o_2,o_1)=\omega_T(o_2,o_1)-h_{o_2\ast o_1}
     +\varphi_{o_1}(h_{o_2})+h_{o_1}
   =\omega_T(o_2,o_1)+(\delta^1 h)(o_2,o_1),
\]
by \eqref{eq:d1}.

\emph{(3) Invariance and the equivalence.} Since the gauge acts transitively on
transversals (Definition~\ref{def:gauge}), (2) shows any two defects are
cohomologous; with (1) this gives a well-defined class $[\omega]\in
H^2(\Sk_\C;\D)$. If \textnormal{(G)} holds then some $T$ has $\omega_T\equiv0$
(Proposition~\ref{prop:closure}), so $[\omega]=0$. Conversely if $[\omega]=0$ then
$\omega_{T}=\delta^1 h$ for some $h$ and any $T$; by (2),
$\omega_{T\cdot(-h)}=\omega_T-\delta^1 h=0$, so $T\cdot(-h)$ closes and
\textnormal{(G)} holds.

\emph{(4) Torsor.} If $T$ closes ($\omega_T=0$), then $T\cdot h$ closes iff
$\omega_{T\cdot h}=\delta^1 h=0$ iff $h\in Z^1$. Thus the closing transversals --
equivalently the strict factorization systems $(\C_T,\D)$ -- form a torsor under
$Z^1(\Sk_\C;\D)$. The $0$-cochains $C^0=\prod_a G_a$ act by
$\delta^0$ as the automorphisms of $(\K,\D)$ that fix every object and rescale the
$\D$-action basepoints; quotienting by $B^0=\operatorname{im}\delta^0$ identifies
SFS related by such an inner relabeling, leaving a torsor under
$Z^1/B^0=H^1(\Sk_\C;\D)$.
\end{proof}

\begin{remark}[Why this is the right complex]
$\Sk_\C$ is the would-be left factor $\C$ \emph{regarded only up to the right
$\D$-action}: its morphisms are the orbits, i.e.\ the $\C$-arrows ``modulo the
ambiguity of choosing a representative.'' The presheaf $\D$ records the residual
torsor of choices. Theorem~\ref{thm:main} says precisely that the local
choices (free, by \textnormal{(L)}) glue to a global section (a closed $\C$) iff a
single $2$-dimensional obstruction vanishes --- the categorical avatar of ``a
principal bundle with flat structure is trivial iff its holonomy class vanishes.''
The companion equivalence ``distributive law $\lambda\colon\D\C\Rightarrow\C\D
\leftrightarrow$ SFS'' of Rosebrugh--Wood \cite{rw} is, in this language, the
statement that the matched-pair cocycle conditions ZS1--ZS4 \cite{pairwise} are the
simplicial identities of which \eqref{eq:d2} is the abelianized shadow.
\end{remark}

\section{The rigid twist computes the generator of $\mathbb Z/2$}

\begin{example}[Ten-morphism rigid twist]\label{ex:rt}
Objects $a,x,y$; $\End(a)=\{\id_a,g\}$ with $g^2=\id_a$; $\End(x),\End(y)$ trivial;
$\Hom(a,x)=\{p,pg\}$, $\Hom(a,y)=\{q,qg\}$, $\Hom(x,y)=\{s,s_2\}$, with the rigid
twist $s\circ p=q$, $s_2\circ p=qg$ (forcing $s\circ pg=qg$, $s_2\circ pg=q$). Take
$\D=\E$ the endomorphisms, so $G_a=\mathbb Z/2$ and $G_x=G_y=0$.
\end{example}

\begin{theorem}\label{thm:rt}
In Example~\ref{ex:rt}, \textnormal{(L)} and \ref{hyp} hold,
$H^2(\Sk_\C;\D)\cong\mathbb Z/2$, and the defect represents the nonzero element;
hence \textnormal{(G)} fails and $\K$ has no strict factorization over $\E$.
\end{theorem}

\begin{proof}
\textnormal{(L)} is the freeness of all six hom-sets over their source vertex
group (verified in \cite{cocycle}). \ref{hyp}(i) holds ($\mathbb Z/2$ abelian).
\ref{hyp}(ii): the only nontrivial $G_b$ is $G_a$, and the only generator into $a$
is $\id_a$ (the regular orbit), on which $g\circ\id_a=\id_a\circ g\in\orb(\id_a)$;
all other targets have trivial $G_b$. So the left action is trivial.

The orbit category $\Sk_\C$ has, besides identities, the morphisms
$[p]\colon a\to x$, $[q]\colon a\to y$, and $[s],[s_2]\colon x\to y$, with
$[s]\ast[p]=\orb(s\circ p)=\orb(q)=[q]=\orb(s_2\circ p)=[s_2]\ast[p]$. The presheaf
is $\D(a)=\mathbb Z/2$, $\D(x)=\D(y)=0$, all restriction maps zero (no nonidentity
generator targets $a$).

Normalized cochains: $C^1$ has the two coordinates $h([p]),h([q])\in\mathbb Z/2$
(generators with source-group $\mathbb Z/2$); $C^2$ has the two coordinates
$\omega([s],[p]),\omega([s_2],[p])\in\mathbb Z/2$ (composable pairs with
$\dom=a$); $C^3=0$ ($y$ is a sink, so no composable triple starts at $a$). Hence
$Z^2=C^2\cong(\mathbb Z/2)^2$. From \eqref{eq:d1} (all $\varphi=0$),
\[
   (\delta^1 h)([s],[p])=h([p])-h([s]\ast[p])=h([p])-h([q]),\quad
   (\delta^1 h)([s_2],[p])=h([p])-h([q]),
\]
so $B^2=\{(t,t):t\in\mathbb Z/2\}$, the diagonal. Therefore
\[
   H^2(\Sk_\C;\D)=(\mathbb Z/2)^2/\langle(1,1)\rangle\cong\mathbb Z/2 .
\]
For the transversal $T=\{p,q,s,s_2,\dots\}$ the defect is
$\omega([s],[p])=0$ (as $s\circ p=q=\flr{}\,\circ\id$) and
$\omega([s_2],[p])=g\mapsto1$ (as $s_2\circ p=qg=q\circ g$), i.e.\
$\omega_T=(0,1)\notin B^2$. Thus $[\omega]\neq0$ is the generator, and by
Theorem~\ref{thm:main}(3), \textnormal{(G)} fails.
\end{proof}

The next proposition is \emph{outside} the hypothesis of Theorem~\ref{thm:main}
(its $\D$ is not a disjoint union of vertex groups), yet exhibits the identical
class --- evidence that the vertex-group restriction in \ref{hyp}(i) is an artifact
of the proof, not of the phenomenon, and can likely be relaxed to ``the defect is
group-valued and the left action is trivial.''

\begin{proposition}[The cross-arrow enlargement keeps the same class]\label{prop:p4}
Adjoin an object $w$ and a free $\langle g\rangle$-torsor $\{d,gd\}\colon w\to a$
to $\D$, with the induced rows $\{pd,pgd\}\colon w\to x$, $\{qd,qgd\}\colon w\to y$
in $\K$ (the pairwise counterexample of \cite{pairwise}). Then $\D$ is \emph{not} a
groupoid ($d$ is noninvertible), so \ref{hyp}(i) fails; nonetheless
\textnormal{(L)} holds, the cross arrows are \emph{absorbed} into existing orbits
($\orb(p)=\{p,pg,pd,pgd\}$ with apex $a$; the only $\C$-arrow into $a$ is still
$\id_a$), the defect remains valued in $G_a=\mathbb Z/2$, the left action is
trivial, and the construction of Theorem~\ref{thm:main} goes through verbatim to
give $H^2(\Sk_\C;\D)\cong\mathbb Z/2$ with the defect the nonzero class.
\end{proposition}

\begin{proof}
Direct computation (\texttt{cohomology\_holonomy.py}): the $\D$-orbit of $p$
contains the $w$-row because $pd=p\circ d$ with $d\in\D$, so $\Sk_\C$ is unchanged
on $a,x,y$ and adds only the source object $w$ (a non-target of any nonidentity
generator). Every defect lands in $\D(a,a)=G_a=\mathbb Z/2$ because each composable
$\C$-pair has the form $(\sigma,c_1)$ with $\cod\sigma=y$, $\dom c_1=a$ and
$\sigma\circ c_1\in\orb(q)$ (apex $a$); the left action is trivial since the only
nontrivial vertex group sits at the \emph{source} $a$, never at a middle object.
The $H^2\cong\mathbb Z/2$ computation is then identical to Theorem~\ref{thm:rt}.
\end{proof}

\section{Honest scope: the nonabelian boundary}

Hypothesis~\ref{hyp}(ii) is not a technicality; it is the exact line past which the
obstruction stops being an abelian $H^2$.

\begin{example}[Nontrivial left action in two morphisms]\label{ex:nonab}
Objects $a,b$; $G_a=0$, $G_b=\{\id_b,h\}$ with $h^2=\id_b$; $\Hom(a,b)=\{t,ht\}$
with $ht:=h\circ t$. Then $\Hom(a,b)$ has two singleton free orbits $\{t\}$,
$\{ht\}$, and the target automorphism $h$ \emph{swaps} them: $h\circ t=ht\notin
\orb(t)$. \ref{hyp}(ii) fails.
\end{example}

When \ref{hyp}(ii) fails, $\orb(c_2\circ c_1)$ depends on the chosen generator of
$o_2$, so $\Sk_\C$ is not a category, the presheaf $\D$ does not exist, and the
coboundary computation in Theorem~\ref{thm:main}(2) acquires a genuinely
\emph{nonabelian} correction: the regauging of $\omega$ mixes the left action
$h_{o_2}\circ c_1$ with the right action, exactly as in a nonabelian $2$-cocycle.
The correct home for the obstruction is then the \emph{nonabelian second
cohomology of the matched pair} $(\C,\D,\lambda)$ --- the categorical analogue of
the Kac/Masuoka classification of bicrossed (matched-pair) extensions by an $H^2$
with the two actions entangled \cite{kac,masuoka}. We state this as the open
continuation:

\begin{quote}\itshape
Conjecture (nonabelian holonomy). For \textnormal{(L)} without \ref{hyp}, the
strict factorizations $(\C,\D)$ of $\K$ extending $\D$ are classified by a
nonabelian $H^2$ of the Rosebrugh--Wood distributive law $\lambda$: the band is the
presheaf-of-groups $a\mapsto G_a$ twisted by the left action, the obstruction to a
$\lambda$ is a class in $H^2$ of that band, and the torsor of $\lambda$'s, when
nonempty, is acted on simply-transitively by $H^1$.
\end{quote}

Example~\ref{ex:nonab} shows the abelian theory is genuinely a special case (there
\textnormal{(G)} happens to hold vacuously, there being no composable $\C$-pairs;
the point is only that $\Sk_\C$ already fails to be a category). The decisive
positive content of this note is that the \emph{entire} single-category rigid-twist
phenomenon --- the seed obstruction, and its crown-jewel reading as the laxator ---
lives in the abelian regime \ref{hyp}, where it is a bona fide $H^2$ class.

\section{Computational verification}\label{sec:comp}

All claims are machine-checked by \texttt{projects/scratch/cohomology\_holonomy.py},
which builds the categories, computes orbits/generators/transversals and the defect
\eqref{eq:defect}, assembles the $\mathbb F_2$ bar complex \eqref{eq:d1}--\eqref{eq:d2},
and locates $[\omega]$.
\begin{itemize}
\item \textbf{Rigid twist.} \textnormal{(L)} holds; no transversal closes;
$\dim_{\mathbb F_2}H^2=1$ ($H^2\cong\mathbb Z/2$); all four transversals give a
defect outside $B^2$ representing the \emph{same} nonzero class
($(1,0)\equiv(0,1)\bmod\langle(1,1)\rangle$). Confirms Theorem~\ref{thm:rt}.
\item \textbf{Pairwise $4$-object.} $\D$ non-groupoid; $\dim H^2=1$; same nonzero
class on all transversals. Confirms Proposition~\ref{prop:p4}.
\item \textbf{Triple chain $a\to x\to y\to z$ (no twist).} \textnormal{(L)} and
\textnormal{(G)} hold; $C^3\neq\varnothing$ and $\delta^2\circ\delta^1=0$ is
verified directly; $\dim H^2=0$. Crucially, six of the eight transversals have
$\omega\neq0$ yet $[\omega]=0$ (they lie in $B^2$) while a closing transversal
\emph{exists} --- demonstrating that the obstruction is the class $[\omega]$, not
the cochain $\omega$. The full additive $2$-cocycle identity is verified on every
composable triple.
\item \textbf{Hypothesis check.} The script verifies abelian vertex groups and the
trivial-left-action condition \ref{hyp}(ii) for the rigid twist and the triple
chain (which satisfy \ref{hyp}(i)) and for the pairwise example (which does not,
but is group-valued with trivial left action --- Proposition~\ref{prop:p4}).
Example~\ref{ex:nonab} is exhibited with $h\circ t\notin\orb(t)$, witnessing that
the trivial-left-action condition genuinely restricts.
\end{itemize}

\section{Grant relevance}

This upgrades \textnormal{(G)} from ``a computable obstruction'' to ``a computable
\emph{invariant}.'' For a compositional-correctness checker the practical payoff is
sharp: after the cheap local freeness test \textnormal{(L)}, one need not search the
whole product of transversal choices for a closed one --- one computes a single
cohomology class $[\omega]\in H^2(\Sk_\C;\D)$, and a strict factorization exists iff
it vanishes; the dimension of $H^2$ measures \emph{how} obstructed the system is,
and $H^1$ enumerates the inequivalent factorizations when they exist. In the
ga-containers framing this realizes the slogan literally: the diversity-collapse
regime $W=1.0$ is the vanishing of a cohomology class. The honest boundary --- that
nontrivial target automorphisms push the obstruction into a nonabelian Kac-type
$H^2$ --- is the precise next theorem, and ties the program to the bicrossed-product
classification literature.

\begin{thebibliography}{9}
\bibitem{pairwise} MacBeth, \emph{The pairwise Zappa--Sz\'ep criterion}, Kodamai
note, 10 June 2026.
\bibitem{cocycle} MacBeth, \emph{The single-category Zappa--Sz\'ep criterion:
explicit cocycle, matched-pair axioms, and why objectwise freeness is necessary but
not sufficient}, Kodamai note, 10 June 2026.
\bibitem{rw} R.~Rosebrugh, R.~J.~Wood, \emph{Distributive laws and factorization},
J.\ Pure Appl.\ Algebra 175 (2002), 327--353.
\bibitem{baueswirsching} H.-J.~Baues, G.~Wirsching, \emph{Cohomology of small
categories}, J.\ Pure Appl.\ Algebra 38 (1985), 187--211.
\bibitem{roos} J.-E.~Roos, \emph{Sur les foncteurs d\'eriv\'es de $\varprojlim$},
C.\ R.\ Acad.\ Sci.\ Paris 252 (1961), 3702--3704.
\bibitem{kac} G.~I.~Kac, \emph{Extensions of groups to ring groups}, Math.\ USSR
Sbornik 5 (1968), 451--474.
\bibitem{masuoka} A.~Masuoka, \emph{Extensions of Hopf algebras and Lie bialgebras},
Trans.\ Amer.\ Math.\ Soc.\ 352 (2000), 3837--3879.
\end{thebibliography}

\end{document}
